% TOP_PROBLEM: {"id": 3207, "title": "Pentagon problem", "category_id": 3, "category": {"id": 3, "name": "graph_theory", "display_name": "Graph Theory", "description": "Problems involving graphs, networks, and their properties.", "slug": "graph-theory", "order_index": 3, "created_at": "2026-07-31T15:26:25.671Z"}, "status": "open", "problem_number": "OPG-167", "set_id": 12}
% FIRST_POSTED: 2026-10-01
% ATTEMPT_STATUS: unresolved
% MODEL: GPT 6 Astra Ultra
% UnsolvedMath ID 3207; source catalogue code OPG-167
% !TeX program = lualatex
% Research attempt dated 2026-10-01; canonical statement preserved.
\documentclass[11pt,a4paper]{article}
\usepackage[margin=25mm]{geometry}
\usepackage{fontspec}
\setmainfont{lmroman10-regular.otf}[BoldFont=lmroman10-bold.otf,
 ItalicFont=lmroman10-italic.otf,BoldItalicFont=lmroman10-bolditalic.otf]
\usepackage{amsmath,amssymb,mathtools}
\usepackage{unicode-math}
\setmathfont{latinmodern-math.otf}
\AtBeginDocument{\let\setminus\smallsetminus}
\usepackage{xurl,hyperref}
\hypersetup{colorlinks=true,urlcolor=blue}
\setcounter{secnumdepth}{0}
\setlength{\parindent}{0pt}
\setlength{\parskip}{5pt}
\setlength{\emergencystretch}{3em}
\begin{document}
\section*{Pentagon problem}\label{3207}
{\small UnsolvedMath ID 3207 · OPG-167 · Graph Theory · OpenGarden}
\subsection{Definitions and mathematical statement}
Let $G$ be a finite simple cubic graph, so each
vertex has exactly three neighbours. Its girth
is the length of its shortest simple cycle.
Let $C_5$ be the cycle with vertex set
$\mathbb Z/5\mathbb Z$, where two residues are
adjacent when their difference is $1$ or $-1$.
A homomorphism $f:G\to C_5$ is a vertex map
preserving adjacency, equivalently a function
satisfying
\[
 f(v)-f(u)\equiv 1\text{ or }-1\pmod5
 \qquad(uv\in E(G)).
\]
The map need not be injective or surjective.

The Pentagon Problem asks whether there
exists an absolute integer $g_0$ such that
every finite simple cubic graph with girth
at least $g_0$ has a homomorphism to $C_5$.
The threshold is chosen once and must work
for all graph orders and all cubic graphs
above it; it cannot depend on $G$.

The target demands more than an ordinary
proper colouring with five colours. Only
consecutive colours around the fixed pentagon
may appear at the ends of an edge.
Equivalently, the vertices are partitioned
into five possibly empty classes, each
independent, with edges permitted only
between consecutive classes cyclically.
The hypothesis forbids all short cycles,
not only short odd cycles. A negative
answer would be a family of cubic graphs
of unbounded girth admitting no pentagon
homomorphism. Disconnected graphs are
allowed and can be mapped component by
component using the same target.
\subsection{Short English statement}
Does one sufficiently large girth threshold guarantee that every cubic graph admits an adjacency-preserving map to a pentagon?
\subsection{Research attempt}
\paragraph{Scope and method.}
This research attempt by GPT-6 Astra Ultra is dated 1 October 2026.
It preserves the catalogue statement and distinguishes elementary proofs
below from cited prior work. No novelty claim or formal proof-assistant
verification is made. The final paragraph states the precise remaining scope.
\paragraph{Literature boundary.}
The problem is attributed to Nešetřil in [S2]. DeVos and Šámal [R1]
prove that cubic graphs of girth at least seventeen map to the Clebsch
graph. Their target is not $C_5$, and their cut-continuous edge mapping
must not be confused with the vertex homomorphism required here.
We investigate the exact consistency equations for a $C_5$ mapping.

\paragraph{A cycle-space formulation.}
Orient every edge arbitrarily. A homomorphism is a potential
$x:V\to\mathbb Z/5\mathbb Z$ whose difference on each oriented edge
belongs to $\{1,-1\}$. Thus it assigns a sign $\epsilon_e$ to every edge,
and on every cycle the signed sum, using traversal direction, is zero
modulo five. These conditions are also sufficient. Choose a root in each
component, set its potential to zero, and integrate the signs along a
spanning tree. For a non-tree edge the fundamental-cycle equation gives
the required potential difference. Consequently it suffices to impose
one equation for each non-tree edge, rather than enumerate all cycles.
This is an exact reduction, but the unknown signs occur simultaneously
in several equations and cannot be chosen separately for each cycle.

\paragraph{Solving individual cycles.}
A closed walk of length $m$ in $C_5$ exists precisely when one can choose
$m$ signs with total divisible by five. Every even $m$ works by taking
equal numbers of positive and negative steps. For odd $m\ge5$, take
$(m+5)/2$ positive and $(m-5)/2$ negative steps, giving total five.
For a triangle, the only totals are $-3,-1,1,3$, none divisible by five.
It follows that every cycle of length at least four maps to $C_5$.
A forest maps by its bipartition onto one fixed target edge.
More generally, a cactus with no triangle maps to $C_5$: attach its edge
and cycle blocks one at a time at their common cutvertex and rotate the
target labels so that the attaching vertex agrees. This proves a useful
compatibility case, but finite cubic graphs need not have a cactus block
structure; the result is not a high-girth cubic theorem.

\paragraph{A necessary global inequality.}
Write $V_i=x^{-1}(i)$ for a proposed mapping. For each $i$ modulo five,
$V_i\cup V_{i+2}$ is independent because those target vertices are
nonadjacent. Each source vertex belongs to exactly two of these five
independent sets. Hence
\[
 5\alpha(G)\ \ge\ \sum_{i=0}^4|V_i\cup V_{i+2}|=2|V(G)|.
\]
Thus $\alpha(G)\ge2|V(G)|/5$ is necessary. This counts actual preimages;
it does not assume the five classes have equal size. It is not sufficient:
a triangle together with enough isolated vertices satisfies the numerical
inequality while still having no mapping. That example is deliberately
outside the cubic high-girth hypothesis and diagnoses only the loss of
information in the independence-number relaxation.

\paragraph{Why a factor-by-factor construction stops.}
Suppose a cubic graph has a perfect matching and all cycles of its
complement are long enough to map. Choose one mapping on each cycle.
Rotating and reversing it gives parameters $a_C\in\mathbb Z/5\mathbb Z$
and $s_C\in\{1,-1\}$. Every matching edge $uv$ between cycles $C,D$
then demands
$a_C+s_Cx(u)-a_D-s_Dx(v)\in\{1,-1\}$.
Local cycle solvability supplies no theorem that this coupled constraint
system is satisfiable. Moreover, restricting to rotations of one chosen
cycle mapping discards other possible sign patterns, so failure of this
restricted search would not disprove the conjecture.

\paragraph{Verification and final status.}
The accompanying root batch script enumerates possible walk residues for
lengths up to forty, checks the cycle criterion, and validates explicit
cubic prism mappings for base cycles of length at least four. Prisms have
four-cycles, so these certificates do not supply an unbounded-girth family.
\textbf{UNRESOLVED.} The exact cycle-space reduction and the necessary
independence bound are proved. A uniform argument satisfying all overlapping
cycle or matching constraints at sufficiently large girth remains missing.

\subsection{Sources}
\begin{itemize}
\item[S1] ULAM.AI and the UnsolvedMath Contributors, supplied problems.json, record 3207 (OPG-167), OpenGarden collection. Catalogue identity and source transcription; CC BY 4.0.
\url{https://huggingface.co/datasets/ulamai/UnsolvedMath}.
\item[S2] Open Problem Garden, Pentagon problem. Original problem and discussion; the mathematical formulation below is expanded from the supplied source record.
\url{http://www.openproblemgarden.org/op/pentagon_problem}.
\item[R1] Matt DeVos and Robert Šámal, \emph{High-girth cubic graphs
are homomorphic to the Clebsch graph}, arXiv:math/0602580.
Primary abstract checked 1 October 2026.
\url{https://arxiv.org/abs/math/0602580}.
\end{itemize}
\end{document}
